\documentclass[11pt]{article}
\usepackage[margin=1in]{geometry}
\usepackage[T1]{fontenc}
\usepackage{lmodern}
\usepackage{amsmath,amssymb,amsthm,mathtools}
\usepackage{booktabs,array}
\usepackage{microtype}
\usepackage{xcolor}
\usepackage[colorlinks=true,linkcolor=blue!45!black,urlcolor=blue!45!black,citecolor=blue!45!black]{hyperref}
\usepackage{enumitem}
\setlist{itemsep=3pt,topsep=5pt}
\setlength{\emergencystretch}{2em}
\allowdisplaybreaks[2]

\newtheorem{theorem}{Theorem}[section]
\newtheorem{lemma}[theorem]{Lemma}
\newtheorem{proposition}[theorem]{Proposition}
\newtheorem{corollary}[theorem]{Corollary}
\theoremstyle{remark}
\newtheorem{remark}[theorem]{Remark}
\numberwithin{equation}{section}

\newcommand{\R}{\mathbb R}
\newcommand{\Z}{\mathbb Z}
\newcommand{\WB}{W_B}
\newcommand{\Ph}{\Phi}
\newcommand{\dd}{\mathrm d}
\newcommand{\cD}{\mathcal D}
\newcommand{\cB}{\mathcal B}
\newcommand{\cH}{\mathcal H}

\title{Uniform Feasibility of the Angular Racah Dual}
\author{}
\date{November 27, 2026}

\begin{document}
\maketitle
\vspace{-1.5em}

\begin{abstract}
We prove feasibility of the terminating Racah dual with degree
$n=\lfloor g/2\rfloor$ on both angular helicity ladders, for every real gap
$g\ge4$, every allowed integer spin, and every real dimension above the twist
gap. The helicity-zero proof uses a contracting quadratic energy for a
variable-coefficient Legendre recurrence. The helicity-four proof uses a
discrete Sonin energy, reducing the entire infinite spin tail to a fixed
degree-five inequality. That inequality has an exact integer certificate:
two explicitly defined bivariate polynomials have 169 positive coefficients
each. Low spins are treated by elementary identities and one displayed
sum-of-squares decomposition. No asymptotic approximation, spin cutoff,
twist discretization, or floating-point sign test enters the proof. The
result gives the angular linear-programming objective
$L=\WB n^2(n^2-1)/400$ and its corresponding exact conditional bound on
$\tau$.
\end{abstract}

\section{The theorem and its normalization}
\label{sec:theorem}

Let $g\ge4$ be real and set
\begin{equation}
 n=\left\lfloor\frac g2\right\rfloor,\qquad
 \nu=n^2,\qquad
 C=J(J+1),\qquad T=(\Delta-1)(\Delta-2).
 \label{eq:basic}
\end{equation}
The heavy-state domains are
\begin{align}
 \cH_0(g)&=\{(J,\Delta):J\in2\Z_{\ge0},\ \Delta\in\R,
                  \ \Delta\ge J+g\},\notag\\
 \cH_4(g)&=\{(J,\Delta):J\in\Z_{\ge4},\ \Delta\in\R,
                  \ \Delta\ge J+g\}.
 \label{eq:domains}
\end{align}
For $b\in\{0,4\}$, define the ladder products
\begin{equation}
 U_k^{(b)}(J,\Delta)
 =\prod_{i=0}^{k-1}
   \frac{C-(b+i)(b+i+1)}{T-(b+i)(b+i+1)},
 \qquad U_0^{(b)}=1,
 \label{eq:ladderU}
\end{equation}
and the terminating polynomial
\begin{equation}
 R_{n,b}(J,\Delta)
 =\sum_{k=0}^{\min(n,J-b)}
    \frac{(-n)_k(n)_k}{(k!)^2}U_k^{(b)}(J,\Delta).
 \label{eq:Rdef}
\end{equation}
Here $(a)_k=a(a+1)\cdots(a+k-1)$, with $(a)_0=1$.
Every denominator in the displayed finite sum is positive on
\eqref{eq:domains}. The weights are
\begin{equation}
 W_b(\Delta)=\frac{T-b(b-1)}3+\frac32,
 \qquad
 \WB=W_4(g+4)=\frac{g^2+5g-3/2}{3}>0.
 \label{eq:weights}
\end{equation}

\begin{theorem}[Uniform heavy feasibility]
\label{thm:main}
For every real $g\ge4$ and $n=\lfloor g/2\rfloor$,
\begin{align}
 W_0(\Delta)+\WB R_{n,0}(J,\Delta)&>0
       &&\text{for all }(J,\Delta)\in\cH_0(g),\label{eq:main0}\\
 W_4(\Delta)-\WB R_{n,4}(J,\Delta)&\ge0
       &&\text{for all }(J,\Delta)\in\cH_4(g).
       \label{eq:main4}
\end{align}
In particular, the coefficients
\begin{equation}
 \boxed{\quad
 \lambda_kD_k^2=-\WB\frac{(-n)_k(n)_k}{(k!)^2},
 \qquad D_k=\binom{k+3}{3},\qquad 0\le k\le n,
 \quad}
 \label{eq:dual}
\end{equation}
are feasible for both heavy ladders in the angular dual linear program.
The helicity-four condition is saturated at $(J,\Delta)=(4,g+4)$.
\end{theorem}

The normalization in \eqref{eq:dual} is the one in the repository's
\texttt{lp\_primal.py}, \texttt{lp\_dual.py}, and
\texttt{racah\_explore.py}, pinned to commit
\texttt{01b01e6811825e446632d23ee8b33aab15b52b59}; see
Section~\ref{sec:reproducibility}. In that normalization,
$r_k^{(b)}=D_k^2U_k^{(b)}$, so
\begin{equation}
 \sum_{k=0}^n\lambda_kr_k^{(b)}=-\WB R_{n,b}.
 \label{eq:normalization}
\end{equation}
The theorem is a statement about the heavy conditions themselves. The precise
linear-programming consequence, with its moment and budget hypotheses stated
explicitly, is given in Section~\ref{sec:LP}.

The proof has two different energy estimates. On the helicity-zero ladder,
a bound slightly weaker than $|R_{n,0}|\le1$ is sufficient because the weight
ratio supplies a strict margin. On the helicity-four ladder, the low-spin
polynomials are bounded above by one. For the remaining spins, a Sonin energy
gives an absolute bound by the actual weight ratio. In particular, the proof
does not require $|R_{n,4}|\le1$, an inequality that fails in general.

\section{An algebraic Racah recurrence at real dimension}
\label{sec:recurrence}

For either helicity, introduce
\begin{equation}
 x=J-b,\qquad N=\Delta-b-2,\qquad s=2b+1,
 \qquad X_x=x(x+s),\qquad Y=N(N+s).
 \label{eq:shifted}
\end{equation}
Then
\begin{equation}
 C-b(b+1)=X_x,\qquad T-b(b+1)=Y,
 \qquad N-x=\Delta-J-2\ge g-2,
 \label{eq:shiftedphysical}
\end{equation}
and
\begin{equation}
 R_{n,b}(J,\Delta)
 =\Ph_x^{(s)}(N)
 :={}_4F_3\left(
 \begin{matrix}-n,n,-x,x+s\\1,-N,N+s\end{matrix};1\right).
 \label{eq:4F3}
\end{equation}
The hypergeometric notation denotes the finite sum in \eqref{eq:Rdef}.
Auxiliary integer indices between zero and a physical spin will be used
below; these auxiliary values do not need to belong to a physical spectrum.

\begin{lemma}[Quadratic-lattice finite differences]
\label{lem:factorial}
Let $s>0$ and
\begin{equation}
 P_k(X)=\prod_{i=0}^{k-1}\bigl(X-i(i+s)\bigr),\qquad
 p_x^{(s)}=\frac{(x+1)(x+s)}{2x+s+1}.
 \label{eq:Pk}
\end{equation}
For integers $k\ge1$, the identity
\begin{align}
 &p_x^{(s)}\bigl[P_k(X_{x+1})-P_k(X_x)\bigr]
 -p_{x-1}^{(s)}\bigl[P_k(X_x)-P_k(X_{x-1})\bigr]\notag\\
 &\hspace{40mm}=k^2(2x+s)P_{k-1}(X_x)
 \label{eq:Pkidentity}
\end{align}
holds as a rational identity in $x$. At $x=0$ the backward term is set to
zero; equivalently $p_{-1}^{(s)}=0$ there.
\end{lemma}

\begin{proof}
Use
\[
 P_k(X_x)=\prod_{i=0}^{k-1}(x-i)(x+s+i).
\]
Away from the zeros of this product, division gives
\begin{align*}
 \frac{P_k(X_{x+1})}{P_k(X_x)}
 &=\frac{(x+1)(x+s+k)}{(x-k+1)(x+s)},\\
 \frac{P_k(X_{x-1})}{P_k(X_x)}
 &=\frac{(x-k)(x+s-1)}{x(x+s+k-1)},\\
 \frac{P_{k-1}(X_x)}{P_k(X_x)}
 &=\frac1{(x-k+1)(x+s+k-1)}.
\end{align*}
Substituting these ratios into \eqref{eq:Pkidentity} proves it by rational
algebra. Clearing denominators extends the identity to zeros of the products.
The boundary value follows directly from the products.
\end{proof}

\begin{proposition}[Exact spin recurrence]
\label{prop:recurrence}
At fixed $N,s,n$, put $y_x=\Ph_x^{(s)}(N)$ and $\nu=n^2$. Then
\begin{equation}
 p_x^{(s)}(y_{x+1}-y_x)-p_{x-1}^{(s)}(y_x-y_{x-1})+q_xy_x=0,
 \qquad q_x=\frac{(2x+s)\nu}{Y-X_x},
 \label{eq:spinrec}
\end{equation}
where the displayed denominators are nonzero. The initial values are
\begin{equation}
 y_0=1,\qquad y_1=1-\frac{(s+1)\nu}{Y}.
 \label{eq:initial}
\end{equation}
This recurrence holds for real $N$.
\end{proposition}

\begin{proof}
Write $y_x=\sum_{k=0}^n a_kP_k(X_x)$, where
\begin{equation}
 a_0=1,\qquad
 a_{k+1}=-\frac{\nu-k^2}{(k+1)^2[Y-k(k+s)]}\,a_k.
 \label{eq:acoeffs}
\end{equation}
Apply Lemma~\ref{lem:factorial} term by term and use
\begin{equation}
 (Y-X)P_k(X)=[Y-k(k+s)]P_k(X)-P_{k+1}(X).
 \label{eq:Pkshift}
\end{equation}
After multiplying the difference operator by $(Y-X_x)/(2x+s)$, the
coefficient of each $P_k(X_x)$ in the left side of \eqref{eq:spinrec}
vanishes by \eqref{eq:acoeffs}. At the terminating endpoint,
$\nu-n^2=0$. The initial values follow from the finite sums at $x=0,1$.
All operations are identities between rational functions; no orthogonality
grid or integer restriction on $N$ is used.
\end{proof}

\section{The helicity-zero ladder}
\label{sec:b0}

Here $b=0$, $s=1$, $Y=T$, and write $F_x=\Ph_x^{(1)}$. The recurrence
\eqref{eq:spinrec} becomes
\begin{equation}
 (x+1)F_{x+1}+xF_{x-1}=(2x+1)c_xF_x,
 \qquad c_x=1-\frac{2\nu}{T-x(x+1)},
 \label{eq:b0rec}
\end{equation}
with $F_0=1$ and $F_1=c_0$.

\subsection{A contracting energy}

For $-1<c<1$, define
\begin{equation}
 \mathcal E_c(u,v)=\frac{u^2+v^2-2cuv}{1-c}
 =(1+c)v^2+\frac{(u-cv)^2}{1-c}.
 \label{eq:E0}
\end{equation}
If $0\le r\le1$ and $w=(1+r)cv-ru$, symmetry in the two arguments gives
\begin{equation}
 \mathcal E_c(v,w)
 =(1+c)v^2+\frac{r^2(u-cv)^2}{1-c}
 \le\mathcal E_c(u,v).
 \label{eq:E0contract}
\end{equation}
Also,
\begin{equation}
 \frac{\partial}{\partial c}\mathcal E_c(u,v)
 =\frac{(u-v)^2}{(1-c)^2}\ge0.
 \label{eq:E0parameter}
\end{equation}
Thus decreasing $c$ decreases this energy.

Suppose $T-J(J+1)>\nu$. Then $-1<c_x<1$ for $0\le x\le J$,
and $c_x$ decreases with $x$. In \eqref{eq:b0rec}, the transfer from
$(F_{x-1},F_x)$ to $(F_x,F_{x+1})$ has $r=x/(x+1)$. Starting from
\[
 \mathcal E_{c_0}(F_0,F_1)=\mathcal E_{c_0}(1,c_0)=1+c_0,
\]
apply \eqref{eq:E0parameter} before each transfer
\eqref{eq:E0contract}. One final change of parameter from $c_{J-1}$
to $c_J$ gives
\begin{equation}
 (1+c_J)F_J^2\le\mathcal E_{c_J}(F_{J-1},F_J)\le1+c_0.
 \label{eq:E0chain}
\end{equation}
Consequently,
\begin{equation}
 \boxed{\quad
 F_J^2\le1+\frac{\nu J(J+1)}{T[T-J(J+1)-\nu]}.
 \quad}
 \label{eq:b0bound}
\end{equation}

\subsection{The complete tail \texorpdfstring{$J\ge4$}{J>=4}}

For $J\ge4$, $g\ge4$, and $\Delta\ge J+g$,
\begin{align}
 T-C
 &\ge(J+g-1)(J+g-2)-J(J+1)\notag\\
 &=(2g-4)J+(g-1)(g-2)\notag\\
 &\ge g^2+5g-14\ge g^2\ge4\nu.
 \label{eq:b0gap}
\end{align}
All denominators in the energy argument are therefore positive, and
\eqref{eq:b0bound} implies
\begin{equation}
 F_J^2\le1+\frac{C}{3T}.
 \label{eq:b0simple}
\end{equation}

Put $B_g=3\WB=g^2+5g-3/2$ and $h=T+9/2-B_g$.
At $T=(J+g-1)(J+g-2)$, direct algebra gives
\begin{align}
 T-B_g&=(J-4)(J+2g+1)+\frac{15}{2}>0,\label{eq:weightslack1}\\
 2h-C&=(J-4)(J+4g-3)+4>0.\label{eq:weightslack2}
\end{align}
Both left sides increase with $T$. Hence $0<B_g<T$ and $h>C/2$, so
\begin{equation}
 \left(\frac{W_0}{\WB}\right)^2
 =\left(1+\frac h{B_g}\right)^2
 >1+\frac C{T}>1+\frac C{3T}\ge F_J^2.
 \label{eq:b0weighted}
\end{equation}
This proves \eqref{eq:main0} for every even $J\ge4$.
For example, the same estimates give the explicit positive margin
\begin{align}
 W_0+\WB F_J
 &\ge W_0-\WB\sqrt{1+C/(3T)}\notag\\
 &\ge\frac h3-\frac{B_gC}{18T}
 >\frac C9.
 \label{eq:b0margin}
\end{align}

\subsection{The spins \texorpdfstring{$J=0,2$}{J=0,2}}

At $J=0$, $F_0=1$, so the condition is immediate. At $J=2$,
\begin{equation}
 F_2=1-\frac{6\nu}{T}+\frac{6\nu(\nu-1)}{T(T-2)}.
 \label{eq:b0spin2}
\end{equation}
The heavy condition gives $T\ge g(g+1)>g^2\ge4\nu$.
Since $n\ge2$, the last term in \eqref{eq:b0spin2} is nonnegative,
and $F_2\ge-1/2$. Moreover,
\begin{equation}
 W_0-\frac{\WB}{2}
 \ge\frac{g^2-3g+21/2}{6}>0.
 \label{eq:b0spin2weight}
\end{equation}
The quadratic in the numerator has positive leading coefficient and negative
discriminant. This completes the entire helicity-zero ladder.

\section{The first five helicity-four spins}
\label{sec:b4low}

For the helicity-four ladder, set
\begin{equation}
 x=J-4,\qquad N=\Delta-6,\qquad Y=N(N+9)=T-20,
 \qquad y_x=\Ph_x^{(9)}(N).
 \label{eq:b4variables}
\end{equation}
Since $W_4(\Delta)\ge\WB$ for $\Delta\ge J+g$ and $J\ge4$,
the one-sided bound $y_x\le1$ suffices at any fixed low spin.

At $x=0$, $y_0=1$. At $x=1$,
\begin{equation}
 1-y_1=\frac{10\nu}{T-20}>0.
 \label{eq:b4x1}
\end{equation}
At $x=2$,
\begin{equation}
 1-y_2=\frac{22\nu(T-3\nu-27)}{(T-20)(T-30)}.
 \label{eq:b4x2}
\end{equation}
Here $T\ge(2n+5)(2n+4)$, and therefore
\[
 T-3n^2-27\ge n^2+18n-7>0.
\]
The denominator in \eqref{eq:b4x2} is positive.

At $x=3$, define
\begin{equation}
 B_3(\nu,T)=36(T-30)(T-42)-234(\nu-1)(T-42)
                  +364(\nu-1)(\nu-4).
 \label{eq:B3}
\end{equation}
Then
\begin{equation}
 1-y_3=\frac{\nu B_3(\nu,T)}{(T-20)(T-30)(T-42)}.
 \label{eq:b4x3}
\end{equation}
With $T_0=(2n+6)(2n+5)$, exact evaluation gives
\begin{align}
 B_3(n^2,T_0)&=4(n^4+297n^3+4405n^2-1089n-338)>0,
 \label{eq:B3initial}\\
 \partial_T B_3(n^2,T_0)&=18(3n^2+88n-11)>0,
 \qquad\partial_T^2 B_3=72.
 \label{eq:B3derivative}
\end{align}
For the first sign, $4405n^2-1089n-338\ge2978n^2$ when $n\ge1$.
It follows that $B_3(n^2,T)>0$ for all $T\ge T_0$, as required.

At $x=4$, put $D_4=Y(Y-10)(Y-22)(Y-36)$. The exact formula is
\begin{equation}
 1-y_4=\frac{26\nu Q(Y,\nu)}{D_4},
 \label{eq:b4x4}
\end{equation}
where
\begin{align}
 Q(Y,\nu)={}&2Y^3-21Y^2\nu-115Y^2+70Y\nu^2+868Y\nu+1806Y\notag\\
            &-70\nu^3-1540\nu^2-7462\nu-6768.
 \label{eq:Q4}
\end{align}
The physical condition gives $N\ge2n+2$. Write $n=m+2$ and
$N=2n+2+t$, with $m,t\ge0$. The sum-of-squares identity in
Appendix~\ref{app:x4} proves $Q>0$ throughout this domain. Since $D_4>0$,
the spin $J=8$ is also feasible. Thus every $x=0,1,2,3,4$ has been covered
for all real dimensions in its heavy domain.

\section{A Sonin energy for the complete helicity-four tail}
\label{sec:sonin}

Keep $N,n$ fixed and use the auxiliary integer index $j$. The recurrence is
\begin{equation}
 p_j\delta_{j+1}-p_{j-1}\delta_j+q_jy_j=0,
 \quad
 p_j=\frac{(j+1)(j+9)}{2j+10},\quad
 q_j=\frac{(2j+9)\nu}{Y-j(j+9)},\quad
 \delta_j=y_j-y_{j-1}.
 \label{eq:b4rec}
\end{equation}

\subsection{Construction and exact decrement}

Define a positive auxiliary sequence by
\begin{equation}
 h_j=2\int_0^1\frac{t^j(1+t^8)}{1+t}\,\dd t.
 \label{eq:hj}
\end{equation}
Its two essential properties are
\begin{equation}
 h_j>h_{j+1}>0,\qquad
 h_j+h_{j+1}=\frac2{j+1}+\frac2{j+9}=\frac2{p_j}.
 \label{eq:hidentity}
\end{equation}
The first follows from $0<t<1$, and the second follows by adding the
integrands. For $j\ge1$, set
\begin{equation}
 E_j=y_j^2+\frac{h_jp_{j-1}p_j}{q_j}\delta_j\delta_{j+1}.
 \label{eq:sonin}
\end{equation}

\begin{lemma}[Exact Sonin identity]
\label{lem:sonin}
Whenever the denominators are nonzero,
\begin{equation}
 E_{j+1}-E_j=
 \left[
 p_j^2\left(\frac{h_{j+1}}{q_{j+1}}-\frac{h_j}{q_j}\right)
 +\frac{p_j}{2}(h_j-h_{j+1})
 \right]\delta_{j+1}^2.
 \label{eq:sonindecrement}
\end{equation}
Also,
\begin{equation}
 E_j=\left(1-\frac{h_jq_j}{4}\right)y_j^2
 +\frac{h_jp_{j-1}^2}{q_j}
       \left(\delta_j-\frac{q_jy_j}{2p_{j-1}}\right)^2.
 \label{eq:soninsquare}
\end{equation}
\end{lemma}

\begin{proof}
In $E_{j+1}-E_j$, eliminate $\delta_{j+2}$ and $\delta_j$ using
\begin{align*}
 \delta_{j+2}&=\frac{p_j\delta_{j+1}-q_{j+1}(y_j+\delta_{j+1})}{p_{j+1}},\\
 \delta_j&=\frac{p_j\delta_{j+1}+q_jy_j}{p_{j-1}}.
\end{align*}
The mixed term in $y_j\delta_{j+1}$ vanishes because
$h_j+h_{j+1}=2/p_j$. The coefficient of $\delta_{j+1}^2$ is the bracket
in \eqref{eq:sonindecrement}. For \eqref{eq:soninsquare}, eliminate
$\delta_{j+1}$ in \eqref{eq:sonin} and complete the square.
\end{proof}

\subsection{Monotonicity throughout the physical interval}

Let $x=J-4$ be the physical endpoint. For every $0\le j\le x$,
\begin{align}
 Y-j(j+9)&=(N-j)(N+j+9)\notag\\
 &\ge(2n-2)(2j+2n+7)\ge4n^2.
 \label{eq:b4gap}
\end{align}
The last inequality follows from
\[
 (2n-2)(2j+2n+7)-4n^2=(4j+10)n-4j-14\ge0
 \qquad(n\ge2).
\]
Consequently,
\begin{equation}
 0<z_j:=\frac{\nu}{Y-j(j+9)}\le\frac14.
 \label{eq:zbound}
\end{equation}
For $j\ge1$,
\begin{equation}
 \frac{p_j}{2j+9}-\frac18
 =\frac{2j^2+21j-9}{8(j+5)(2j+9)}\ge0.
 \label{eq:pbound}
\end{equation}
Equations \eqref{eq:zbound}--\eqref{eq:pbound} imply $q_j\le2p_j$.
The sequence $q_j$ is increasing as long as its denominator is positive.
Hence the bracket in \eqref{eq:sonindecrement} is at most
\begin{equation}
 (h_j-h_{j+1})p_j\left(\frac12-\frac{p_j}{q_j}\right)\le0.
 \label{eq:decrementsign}
\end{equation}
We have therefore proved $E_{j+1}\le E_j$ for every step needed to reach
the physical endpoint.

\subsection{A uniform lower bound for the energy}

For $0\le t\le1$, the inequality $1/(1+t)\le1-t/2$ gives
\begin{equation}
 h_j\le H_j:=\frac2{j+1}+\frac2{j+9}-\frac1{j+2}-\frac1{j+10}.
 \label{eq:hupper}
\end{equation}
For $j\ge5$,
\begin{equation}
 \frac43-\frac{2j+9}{4}H_j
 =\frac{2j^4+41j^3+160j^2-869j-2502}
 {6(j+1)(j+2)(j+9)(j+10)}\ge0.
 \label{eq:coercivitypoly}
\end{equation}
Indeed, after $j=5+v$, the numerator equals
$2v^4+81v^3+1075v^2+4806v+3528$. Together with \eqref{eq:zbound}, this
implies $h_jq_j/4\le1/3$. The square completion
\eqref{eq:soninsquare} therefore gives
\begin{equation}
 \boxed{\qquad E_j\ge\frac23y_j^2\qquad(5\le j\le x).\qquad}
 \label{eq:coercivity}
\end{equation}
Combining monotonicity with \eqref{eq:coercivity},
\begin{equation}
 \boxed{\qquad y_x^2\le\frac32E_x\le\frac32E_5\qquad(x\ge5).\qquad}
 \label{eq:tailreduction}
\end{equation}

All auxiliary denominators are positive because $N-j\ge g-2\ge2$ for
$j\le x$. The expression $E_x$ contains $y_{x+1}$, which is well-defined
by the finite sum and the recurrence. It is only an auxiliary algebraic
value. The decrement identity is used for $j=5,\ldots,x-1$, and requires
no physical-state assumption at $x+1$.

\section{The exact degree-five certificate}
\label{sec:certificate}

The infinite tail has now been reduced to a fixed calculation. The following
inequality is an exact polynomial statement over a real continuous domain.

\begin{lemma}[Fixed-energy bound]
\label{lem:E5}
Let $n\ge1$, $\nu=n^2$, and $Y\ge4n^2+30n+36$. Let $y_4,y_5$ be
the finite polynomials in \eqref{eq:4F3} with $s=9$, expressed as rational
functions of $\nu,Y$. For $0\le h\le4/15$, define
\begin{equation}
 \mathcal E_5(h)=y_5^2
 +h\frac{4225(Y-70)}{6156\nu}(y_5-y_4)^2
 -h\frac{65}{18}y_5(y_5-y_4).
 \label{eq:E5direct}
\end{equation}
Then
\begin{equation}
 \mathcal E_5(h)
 <\frac23\left(\frac{Y+25/2}{4n^2+18n+25/2}\right)^2.
 \label{eq:E5certificate}
\end{equation}
In particular, this applies to $E_5=\mathcal E_5(h_5)$.
\end{lemma}

The proof is given explicitly below. It combines algebraic identities with
a finite integer-coefficient certificate. The full coefficients are recorded
in \texttt{coefficient\_ledger.json}; the accompanying
\texttt{verify\_angular\_racah\_feasibility.py} reconstructs the polynomials
and checks every coefficient with exact arithmetic.

\subsection{Explicit polynomials}

Define
\begin{equation}
 \cD=Y(Y-10)(Y-22)(Y-36)(Y-52),\qquad
 A=\cD y_5,\qquad P=\frac{\cD(y_4-y_5)}{18\nu}.
 \label{eq:APdefs}
\end{equation}
These are polynomials in $\nu,Y$. A factored expression for $A$ is
\begin{align}
 A={}&\cD-70\nu(Y-10)(Y-22)(Y-36)(Y-52)\notag\\
 &+1050\nu(\nu-1)(Y-22)(Y-36)(Y-52)\notag\\
 &-5600\nu(\nu-1)(\nu-4)(Y-36)(Y-52)\notag\\
 &+11900\nu(\nu-1)(\nu-4)(\nu-9)(Y-52)\notag\\
 &-8568\nu(\nu-1)(\nu-4)(\nu-9)(\nu-16).
 \label{eq:Aexplicit}
\end{align}
The second polynomial is
\begin{align}
 P={}&Y^4-28Y^3\nu-92Y^3
       +210Y^2\nu^2+2030Y^2\nu+2668Y^2\notag\\
 &-560Y\nu^3-10640Y\nu^2-41664Y\nu-26400Y\notag\\
 &+476\nu^4+14840\nu^3+115388\nu^2+224112\nu+57024.
 \label{eq:Pexplicit}
\end{align}
Set
\begin{equation}
 \cB=4225\nu(Y-70)P^2+1235\nu AP.
 \label{eq:Bexplicit}
\end{equation}
Substituting \eqref{eq:APdefs} into \eqref{eq:E5direct} gives the exact
identity
\begin{equation}
 \mathcal E_5(h)=\frac{19A^2+h\cB}{19\cD^2}.
 \label{eq:E5rational}
\end{equation}
The cancellation of $\nu$ in \eqref{eq:APdefs} follows either directly
from the finite sums or from $y_4-y_5=0$ at $\nu=0$.

\subsection{Two integer polynomials}

For $H=0,4$, define
\begin{equation}
 F_H(n,Y)=190(2Y+25)^2\cD^2
 -(8n^2+36n+25)^2\bigl(285A^2+H\cB\bigr),
 \qquad \nu=n^2.
 \label{eq:FH}
\end{equation}
Equivalently,
\begin{align}
 F_H=60\biggl[&\frac23(Y+25/2)^2(19\cD^2)\notag\\
 &-(4n^2+18n+25/2)^2
       \left(19A^2+\frac H{15}\cB\right)\biggr].
 \label{eq:FHequiv}
\end{align}
Now substitute
\begin{equation}
 n=u+1,\qquad Y=4n^2+30n+36+t,
 \qquad u,t\ge0.
 \label{eq:certsubstitution}
\end{equation}
Each resulting polynomial has the exact support
\begin{equation}
 F_H=\sum_{b=0}^{12}\sum_{a=0}^{24-2b}c^{(H)}_{a,b}u^at^b,
 \qquad c^{(H)}_{a,b}\in\Z_{>0}.
 \label{eq:certificatepositive}
\end{equation}
There are $\sum_{b=0}^{12}(25-2b)=169$ coefficients for each value of
$H$. Table~\ref{tab:coefficients} lists the smallest coefficient at each
fixed power $b$ of $t$. In every row this minimum occurs at $a=24-2b$.
All coefficients, including these minima, are obtained by expansion of the
explicit integer expressions \eqref{eq:Aexplicit}--\eqref{eq:FH} after
\eqref{eq:certsubstitution}.

\begin{table}[htbp]
\centering
\caption{Exact coefficient minima in the two degree-five energy certificates.
For each row, $\min_{0\le a\le24-2b}c^{(H)}_{a,b}$ is shown.}
\label{tab:coefficients}
\renewcommand{\arraystretch}{1.12}
\begin{tabular}{@{}rrr@{}}
\toprule
$b$ & $H=0$ & $H=4$\\
\midrule
0 & 11479429120 & 912363520\\
1 & 46534471680 & 12384399360\\
2 & 47966530560 & 26301250560\\
3 & 15834265600 & 29297489920\\
4 & 5911219200 & 20565335040\\
5 & 8419584000 & 9314150400\\
6 & 4375338240 & 2730174720\\
7 & 746964480 & 551278080\\
8 & 47424000 & 88657920\\
9 & 12524800 & 11127040\\
10 & 784320 & 784320\\
11 & 36480 & 36480\\
12 & 760 & 760\\
\bottomrule
\end{tabular}
\end{table}

\begin{proof}[Proof of Lemma~\ref{lem:E5}]
The coefficient identity \eqref{eq:certificatepositive} shows that $F_0$
and $F_4$ are strictly positive for all $u,t\ge0$, including $u=t=0$.
Because \eqref{eq:FH} is affine in $H$, it is strictly positive for every
$H\in[0,4]$. Thus \eqref{eq:FHequiv} proves
\eqref{eq:E5certificate} for every $h=H/15\in[0,4/15]$, after division
by the positive factors $60$, $19\cD^2$, and
$(4n^2+18n+25/2)^2$.

There is no zero denominator on the stated domain: $Y\ge70>52$ for
$n\ge1$. Formula \eqref{eq:E5direct} remains regular at $Y=70$.
For the physical theorem, $n\ge2$ gives the stronger $Y\ge112$.
Finally, \eqref{eq:hupper} at $j=5$ gives
\begin{equation}
 0<h_5\le\frac26+\frac2{14}-\frac17-\frac1{15}=\frac4{15}.
 \label{eq:h5interval}
\end{equation}
The recurrence \eqref{eq:b4rec} eliminates $y_6$ from the Sonin definition
and gives exactly \eqref{eq:E5direct} with $h=h_5$:
$p_4=65/18$ and $q_5=19\nu/(Y-70)$. This proves the asserted estimate
for the actual energy $E_5$.
\end{proof}

The integer expansion is the finite computational component of the proof.
Its input is a fully specified pair of polynomials, and its output is their
complete coefficient list. Positivity of those coefficients proves an
inequality on the full real quadrant $u,t\ge0$. Thus this step supplies a
continuum certificate, with no sampling parameter or numerical precision
assumption.

\section{Completion of the heavy-feasibility theorem}
\label{sec:completion}

Only the helicity-four spins $x=J-4\ge5$ remain. The physical condition
and the definition of $n$ give
\begin{equation}
 N\ge x+g-2\ge g+3\ge2n+3,
 \qquad Y\ge(2n+3)(2n+12)=4n^2+30n+36.
 \label{eq:E5physical}
\end{equation}
The exact weights in these variables are
\begin{equation}
 W_4(\Delta)=\frac{Y+25/2}{3},\qquad
 \WB=\frac{g^2+5g-3/2}{3}.
 \label{eq:b4weights}
\end{equation}
Since $g<2n+2$,
\begin{equation}
 g^2+5g-\frac32\le4n^2+18n+\frac{25}{2}.
 \label{eq:floorweight}
\end{equation}
Combining \eqref{eq:tailreduction}, Lemma~\ref{lem:E5}, and
\eqref{eq:floorweight},
\begin{equation}
 y_x^2\le\frac32E_5
 <\left(\frac{Y+25/2}{4n^2+18n+25/2}\right)^2
 \le\left(\frac{W_4(\Delta)}{\WB}\right)^2.
 \label{eq:fulltailbound}
\end{equation}
The weights are positive, so \eqref{eq:main4} holds strictly in this tail.
Section~\ref{sec:b4low} covers $x=0,1,2,3,4$, and
Section~\ref{sec:b0} covers every allowed helicity-zero spin. This proves
Theorem~\ref{thm:main}.

The floor prescription enters the proof in two distinct ways:
$2n\le g$ gives the lower bounds used in the recurrences, while
$g<2n+2$ gives the weight comparison \eqref{eq:floorweight}.
Both inequalities are retained in the exact theorem. The phrase
``$n$ approximately $g/2$'' is therefore replaced by the precise choice
$n=\lfloor g/2\rfloor$ throughout.

\begin{remark}[The proportionality constant matters]
\label{rem:ratio}
For fixed $x=J-4$ and $\Delta=J+g$, if $n/g\to\rho$ as $g\to\infty$,
the finite sum has the limit
\[
 y_x\longrightarrow{}_2F_1(-x,x+9;1;\rho^2),
 \qquad W_4/\WB\longrightarrow1.
\]
At $x=2$, the necessary limiting inequality is $\rho^2\le1/3$.
At $x=3$, the limit is $1-36\rho^2+234\rho^4-364\rho^6$.
Combining these two low-spin conditions gives the necessary bound
\begin{equation}
 \rho\le\sqrt{\frac{117-\sqrt{585}}{364}}
 =0.504956822\ldots.
 \label{eq:rhomax}
\end{equation}
In particular, the theorem cannot be extended by replacing $1/2$ with
$0.51$ in the proportional degree choice. This observation concerns this
specific Racah family and does not restrict other angular duals.
\end{remark}

\section{The exact angular linear-programming corollary}
\label{sec:LP}

The proven heavy conditions immediately determine the objective value of
this particular dual. Since $D_2=10$ and
\[
 \frac{(-n)_2(n)_2}{(2!)^2}=\frac{n^2(n^2-1)}4,
\]
equation \eqref{eq:dual} gives
\begin{equation}
 \boxed{\quad L=-\lambda_2=\frac{\WB n^2(n^2-1)}{400}>0.\quad}
 \label{eq:L}
\end{equation}

For clarity, the hypotheses under which this gives a bound on $\tau$ can
be stated directly in the repository's primal normalization. Let
$\rho_0,\rho_4$ be nonnegative heavy measures supported on
$\cH_0(g),\cH_4(g)$. Suppose their moments exist and satisfy the angular
harmonic-balance relations through degree $n$,
\begin{equation}
 \int r_k^{(4)}\,\dd\rho_4-\int r_k^{(0)}\,\dd\rho_0
 =\begin{cases}|c|^2,&k=2,\\0,&k\ne2,\end{cases}
 \qquad 0\le k\le n,
 \label{eq:LPmoments}
\end{equation}
and the endpoint budget
\begin{equation}
 \int W_4\,\dd\rho_4+\int W_0\,\dd\rho_0\le3+D,
 \qquad 3+D\ge0.
 \label{eq:LPbudget}
\end{equation}
For a discrete spectrum these integrals are sums. The case $D=0$ is the
endpoint budget used in \texttt{lp\_primal.py}; the displayed form keeps
an explicitly specified budget correction.

\begin{corollary}[Conditional angular LP bound]
\label{cor:LP}
Under \eqref{eq:LPmoments}--\eqref{eq:LPbudget}, with
$\tau=4|c|$, the dual \eqref{eq:dual} gives
\begin{equation}
 |c|^2\le\frac{3+D}{L},\qquad
 \boxed{\quad
 \tau\le\frac{80\sqrt{3+D}}{\sqrt{\WB}\,n\sqrt{n^2-1}}.
 \quad}
 \label{eq:taubound}
\end{equation}
For $D=0$, the right side has the asymptotic expansion
\begin{equation}
 \frac{80\sqrt3}{\sqrt{\WB}\,n\sqrt{n^2-1}}
 =\frac{960}{g^3}\left(1+O(g^{-1})\right),\qquad g\longrightarrow\infty.
 \label{eq:tauasymptotic}
\end{equation}
\end{corollary}

\begin{proof}
Integrate both feasible heavy conditions in their form
\eqref{eq:normalization}. Nonnegativity of the measures and
\eqref{eq:LPmoments} imply
\begin{align*}
 0
 &\le\int\left(W_4+\sum_k\lambda_kr_k^{(4)}\right)\dd\rho_4
   +\int\left(W_0-\sum_k\lambda_kr_k^{(0)}\right)\dd\rho_0\\
 &=\int W_4\,\dd\rho_4+\int W_0\,\dd\rho_0-L|c|^2.
\end{align*}
The budget gives $L|c|^2\le3+D$. Substituting \eqref{eq:L} and
$\tau=4|c|$ proves \eqref{eq:taubound}.
For \eqref{eq:tauasymptotic}, use
\[
 \WB=\frac{g^2}{3}\bigl(1+O(g^{-1})\bigr),\qquad
 n=\frac g2+O(1),\qquad n\sqrt{n^2-1}=\frac{g^2}{4}\bigl(1+O(g^{-1})\bigr).
\]
\end{proof}

The exact output of the theorem is heavy feasibility and the objective
\eqref{eq:L}. Corollary~\ref{cor:LP} applies when the displayed moment
relations and endpoint budget are available. It retains the correction
$D$ explicitly; establishing or estimating that correction in a particular
CFT is a separate input to the corollary.

\section{Reproducibility and source provenance}
\label{sec:reproducibility}

The normalization was read from the directory
\texttt{ads4/code/angular\_opt/} in
\href{https://github.com/anshM123/MAJOR}{\texttt{anshM123/MAJOR}}, at the
immutable commit
\begin{center}
\texttt{01b01e6811825e446632d23ee8b33aab15b52b59}.
\end{center}
The relevant source files are \texttt{lp\_primal.py},
\texttt{lp\_dual.py}, and \texttt{racah\_explore.py}. The proof is
expressed above entirely in finite products, rational recurrences, positive
integrals, and explicit polynomial identities, so it can be checked without
running the source repository's optimization or sampling routines.

The accompanying verifier and coefficient ledger provide the exact finite
arithmetic checks used by the proof:
\begin{itemize}
 \item \texttt{verify\_angular\_racah\_feasibility.py} reconstructs the
 recurrence and energy identities and the fixed finite certificates.
 \item \texttt{coefficient\_ledger.json} records every integer
 $c^{(H)}_{a,b}$ in \eqref{eq:certificatepositive}, for $H=0$ and $H=4$.
 The support and the row minima are specified in
 \eqref{eq:certificatepositive} and Table~\ref{tab:coefficients}.
\end{itemize}
The unified verifier completes 44 exact checks, including all 338 integer
coefficients in the two endpoint certificates; their overall minimum is
$760$ in the integer normalization \eqref{eq:FH}.
The degree-five energy calculation was also independently reconstructed from
the spin recurrence, starting from $y_0=1$ and $y_1=1-10\nu/Y$.
That reconstruction gives the same coefficient signs without using the
finite-sum expansion for $A$ or $P$. Both versions use exact rational
arithmetic. The all-spin and all-real-dimension parts of the argument are
the analytic energy estimates in Sections~\ref{sec:b0} and
\ref{sec:sonin}; the fixed coefficient check supplies their remaining
finite endpoint inequality.

\appendix
\section{The explicit sum-of-squares certificate at \texorpdfstring{$J=8$}{J=8}}
\label{app:x4}

For $x=4$, the finite sum gives \eqref{eq:b4x4} with $Q$ in
\eqref{eq:Q4}. The physical assumptions imply $n\ge2$ and $N\ge2n+2$.
Set
\begin{equation}
 n=m+2,\qquad N=2n+2+t,\qquad m,t\ge0,
 \qquad \nu=n^2,\quad Y=N(N+9).
 \label{eq:x4substitution}
\end{equation}
Then the following exact identity holds:
\begin{align}
 Q={}&2m^4(m-2t-7)^2+m^4(38t^2+462t+1578)\notag\\
 &+m^3(152t^3+3332t^2+23036t+52080)\notag\\
 &+m^2(99t^4+3486t^3+40027t^2+182448t+281098)\notag\\
 &+m(24t^5+1176t^4+20368t^3+150612t^2+480668t+546588)\notag\\
 &+2t^6+126t^5+2987t^4+32844t^3+169559t^2+402738t+356184.
 \label{eq:x4SOS}
\end{align}
Every term is nonnegative for $m,t\ge0$, and the constant term is strictly
positive. Thus $Q>0$. Moreover, $N\ge6$ gives $Y\ge90$, so
$D_4=Y(Y-10)(Y-22)(Y-36)>0$. Equation \eqref{eq:b4x4} proves $y_4<1$
uniformly on the whole heavy region, including noninteger dimensions.

The identity \eqref{eq:x4SOS} can be verified by substituting
\eqref{eq:x4substitution} into \eqref{eq:Q4} and collecting powers of
$m,t$. Its single square absorbs the two negative coefficients in the
direct monomial expansion; every coefficient of the remaining polynomial
is positive.

\end{document}
